\chapter{Vorbemerkungen}

Alle hier behandelten Schaltungen entsprechen der sog. Positiv-Logik. Das heißt, dass der H-Bereich eine positive Spannung sein muss. Es gibt auch eine Negativ-Logik, die aber in der Technik kaum eine Bedeutung hat.\index{Positiv-Logik}\index{H-Bereich}\index{Spannung}\index{Negativ-Logik}

Weiterhin gilt ein „offener“ Eingang (nicht angeschlossener Eingang) immer als Wert 1 (entspr. H). Dies ist durch die Bauweise der ICs (Integrierte Schaltungen) bedingt, die intern immer den Wert 1 an einen Ausgang legen, wenn dieser nicht angeschlossen ist. Das bringt vor allem beim Zusammenschalten mehrerer Verknüpfungsschaltungen enorme Vorteile.\index{Integrierte Schaltungen}\index{Verknüpfungsschaltungen}

Der Masseanschluss bei digitalen Schaltungen hat immer den Wert 0. Dies entspricht bei der Stromversorgung dem Minuspol. Der Wert 1 ist also der Pluspol der Spannungsquelle. Falls Messungen mit einem Spannungsmessgerät oder gar Oszilloskop gemacht werden, ist daher der Massepunkt immer der negative Pol der Spannungsquelle. Die kapazitive Last für digitale Schaltungen sollte 100 pF nicht überschreiten, damit die Funktion gewährleistet ist und eine Überlast der Ausgänge durch zu hohe Auf- bzw. Entladeströme vermieden wird. Sind größere Kondensatoren zur Signalverzögerung erforderlich, so ist ein Vorwiderstand vorzusehen.\index{Masseanschluss}\index{Spannungsquelle}\index{Spannungsmessgerät}\index{Oszilloskop}\index{kapazitive Last}\index{Signalverzögerung}

Der Ausgangsstrom von TTL-Schaltungen ist sehr gering, meist bei 20 mA. Das reicht wohl für den Betrieb einer Kontroll-LED (Licht Emittierende Diode = Leuchtdiode), doch ist darauf zu achten, dass der Ausgang nicht überlastet wird. Notfalls muss ein Schaltverstärker dem Ausgang folgen, der das Ausgangssignal verstärkt. Bei manchen IC’s jedoch ist speziell für LED-Betrieb ein Treiber vorhanden, der einen Anschluss von mehr als einer LED erlaubt (z.B. beim 7-Segment-Decoder). Es ist ratsam, sich deshalb in einem Datenbuch über diesen Sachverhalt zu informieren.\index{Ausgangsstrom}\index{LED}\index{Ausgangssignal}\index{Treiber}

